\section{Schauder fixed point theorem}

\begin{lemma}[Bump functions]
\label{lem:schauder-bumps}
\lean{SchauderFixedPointTheorem.bump_continuous,SchauderFixedPointTheorem.sum_bump_pos}
\leanok
The bump functions used in the finite-dimensional approximation are continuous and their sum is positive on a covered compact set.
\end{lemma}

\begin{lemma}[Center-of-mass approximation]
\label{lem:schauder-center}
\lean{SchauderFixedPointTheorem.centerMass_approx}
\leanok
\uses{lem:schauder-bumps}
The bump-weighted center of mass stays within the prescribed approximation radius.
\end{lemma}

\begin{theorem}[Schauder projection]
\label{thm:schauder-projection}
\lean{SchauderFixedPointTheorem.schauder_projection_exists}
\leanok
\uses{lem:schauder-bumps,lem:schauder-center}
A compact set admits a continuous approximation by the convex hull of finitely many points.
\end{theorem}

\begin{theorem}[Finite-dimensional Brouwer step]
\label{thm:schauder-brouwer}
\lean{SchauderFixedPointTheorem.brouwer_fixed_point_convexHull_finite}
\notready
A continuous self-map of the convex hull of a finite nonempty set has a fixed point. This Brouwer fixed-point input is still a \texttt{sorry}.
\end{theorem}

\begin{theorem}[Schauder fixed point theorem]
\label{thm:schauder-fixed-point}
\lean{SchauderFixedPointTheorem.MainTheorem}
\notready
\uses{thm:schauder-projection,thm:schauder-brouwer}
A continuous self-map of a nonempty compact convex subset of a Banach space has a fixed point.
\end{theorem}
